% --- Archon physics formalization source begin ---
% archon:physics
% archon:covers PhyXMiniProblems/problem_phyx_mini_0896.lean
% archon:source-report reports/phyx_mini/problem_phyx_mini_0896.source.json

\chapter{Physics problem phyx\_mini\_0896}
\label{ch:PhyXMiniProblems_problem_phyx_mini_0896}

\paragraph{Problem source.}
\#\# Question

What is the magnitude of the force \textbackslash{}( \textbackslash{}vec\{F\} \textbackslash{}) on the \textbackslash{}( -10\textbackslash{},\textbackslash{}mathrm\{nC\} \textbackslash{}) charge in figure?

\#\# Answer choices

- A: A: 7.25 \textbackslash{}times 10\textasciicircum{}\{-3\}N

- B: B: 6.35 \textbackslash{}times 10\textasciicircum{}\{-3\}N

- C: C: 7.3 \textbackslash{}times 10\textasciicircum{}\{-3\}N

- D: D: 1.33 \textbackslash{}times 10\textasciicircum{}\{-3\}N

\#\# Figure caption (auxiliary; use the image as primary evidence)

The image depicts a simple diagram involving three point charges placed at the corners of a right triangle. 

- The bottom-left charge is labeled as \textbackslash{}(10 \textbackslash{}, \textbackslash{}text\{nC\}\textbackslash{}) with a positive sign, indicating a positive charge. It is depicted in red.

- The bottom-right charge is labeled as \textbackslash{}(-10 \textbackslash{}, \textbackslash{}text\{nC\}\textbackslash{}) with a negative sign, indicating a negative charge. It is depicted in teal.

- The top-right charge is labeled as \textbackslash{}(8.0 \textbackslash{}, \textbackslash{}text\{nC\}\textbackslash{}) with a positive sign, indicating a positive charge. It is also depicted in red.

- The horizontal distance between the charges labeled \textbackslash{}(10 \textbackslash{}, \textbackslash{}text\{nC\}\textbackslash{}) and \textbackslash{}(8.0 \textbackslash{}, \textbackslash{}text\{nC\}\textbackslash{}) is given as \textbackslash{}(3.0 \textbackslash{}, \textbackslash{}text\{cm\}\textbackslash{}).

- The vertical distance between the charges labeled \textbackslash{}(8.0 \textbackslash{}, \textbackslash{}text\{nC\}\textbackslash{}) and \textbackslash{}(-10 \textbackslash{}, \textbackslash{}text\{nC\}\textbackslash{}) is given as \textbackslash{}(1.0 \textbackslash{}, \textbackslash{}text\{cm\}\textbackslash{}).

The lines connecting the charges are dashed, indicating the geometric relationships between them.

\#\# Dataset metadata

Electrostatics; Spatial Relation Reasoning, Multi-Formula Reasoning

\paragraph{Recorded answer/context.}
C: C: 7.3 \textbackslash{}times 10\textasciicircum{}\{-3\}N

\paragraph{Figure/image path.}
/root/proposal\_for\_physic/science-mango/phyx\_mini\_run/phyx\_data/test\_image/896.png

\paragraph{Formalization target.}
create a compiling Lean file with sorry bodies at `PhyXMiniProblems/problem\_phyx\_mini\_0896.lean`. The Lean declarations must preserve the physical quantities, dimensions or dimensional roles, figure labels, governing-law hypotheses, and final relation expressed by this problem.
Use Mathlib/Physlib names found through LeanExplore where available. If a physics API is missing, introduce faithful local abstractions rather than scalar placeholder aliases.

\begin{theorem}[Physics formalization target]
\label{thm:physics:phyx_mini_0896:target}
\lean{PhyXMiniProblems.ProblemPhyXMini0896.problem_phyx_mini_0896}
\uses{def:physics:phyx-mini-0896:phyxminiproblems-problemphyxmini0896-forcemagnitudeinnewtons, def:physics:phyx-mini-0896:phyxminiproblems-problemphyxmini0896-threepointchargesetup, def:physics:phyx-mini-0896:phyxminiproblems-problemphyxmini0896-matchesquestiontarget, def:physics:phyx-mini-0896:phyxminiproblems-problemphyxmini0896-matchessuppliedthreechargefigure, def:physics:phyx-mini-0896:phyxminiproblems-problemphyxmini0896-usesvacuumcoulombconstant, def:physics:phyx-mini-0896:phyxminiproblems-problemphyxmini0896-hasphysicalthreechargeparameters, def:physics:phyx-mini-0896:phyxminiproblems-problemphyxmini0896-satisfiespointchargecoulombforcelaw, def:physics:phyx-mini-0896:phyxminiproblems-problemphyxmini0896-satisfieselectrostaticforcesuperposition, def:physics:phyx-mini-0896:phyxminiproblems-problemphyxmini0896-recordeddatasetanswer, def:physics:phyx-mini-0896:phyxminiproblems-problemphyxmini0896-roundstonearestquantum}
The assigned autoformalize agent should translate this physics problem into Lean declarations in the covered file, with theorem and lemma proof bodies written as `by sorry`.
\end{theorem}
\begin{proof}
This is an autoformalization task, not a proof task. Produce statements that can later be proved by the physics prover without weakening the physical model.
\end{proof}

% --- Archon declaration topology begin ---
\section{Declaration topology}
\begin{definition}[force Dimension]
  \label{def:physics:phyx-mini-0896:phyxminiproblems-problemphyxmini0896-forcedimension}
  \lean{PhyXMiniProblems.ProblemPhyXMini0896.forceDimension}
  The physical dimension of force, M L T⁻²
\end{definition}
\begin{proof}
  This is the declared model definition of force Dimension.
\end{proof}
\begin{definition}[Signed Charge Quantity]
  \label{def:physics:phyx-mini-0896:phyxminiproblems-problemphyxmini0896-signedchargequantity}
  \lean{PhyXMiniProblems.ProblemPhyXMini0896.SignedChargeQuantity}
  A signed, unit-independent electric charge
\end{definition}
\begin{proof}
  This is the declared model definition of Signed Charge Quantity.
\end{proof}
\begin{definition}[Length Quantity]
  \label{def:physics:phyx-mini-0896:phyxminiproblems-problemphyxmini0896-lengthquantity}
  \lean{PhyXMiniProblems.ProblemPhyXMini0896.LengthQuantity}
  A nonnegative, unit-independent physical length
\end{definition}
\begin{proof}
  This is the declared model definition of Length Quantity.
\end{proof}
\begin{definition}[Planar Force Quantity]
  \label{def:physics:phyx-mini-0896:phyxminiproblems-problemphyxmini0896-planarforcequantity}
  \lean{PhyXMiniProblems.ProblemPhyXMini0896.PlanarForceQuantity}
  \uses{def:physics:phyx-mini-0896:phyxminiproblems-problemphyxmini0896-forcedimension}
  A unit-independent force vector in the plane of the figure
\end{definition}
\begin{proof}
  This is the declared model definition of Planar Force Quantity.
\end{proof}
\begin{definition}[charge In Coulombs]
  \label{def:physics:phyx-mini-0896:phyxminiproblems-problemphyxmini0896-chargeincoulombs}
  \lean{PhyXMiniProblems.ProblemPhyXMini0896.chargeInCoulombs}
  \uses{def:physics:phyx-mini-0896:phyxminiproblems-problemphyxmini0896-signedchargequantity}
  Coherent-SI coulomb readout of a signed physical charge
\end{definition}
\begin{proof}
  This is the declared model definition of charge In Coulombs.
\end{proof}
\begin{definition}[charge In Nanocoulombs]
  \label{def:physics:phyx-mini-0896:phyxminiproblems-problemphyxmini0896-chargeinnanocoulombs}
  \lean{PhyXMiniProblems.ProblemPhyXMini0896.chargeInNanocoulombs}
  \uses{def:physics:phyx-mini-0896:phyxminiproblems-problemphyxmini0896-signedchargequantity, def:physics:phyx-mini-0896:phyxminiproblems-problemphyxmini0896-chargeincoulombs}
  Nanocoulomb readout used by the three printed charge labels
\end{definition}
\begin{proof}
  This is the declared model definition of charge In Nanocoulombs.
\end{proof}
\begin{definition}[length In Meters]
  \label{def:physics:phyx-mini-0896:phyxminiproblems-problemphyxmini0896-lengthinmeters}
  \lean{PhyXMiniProblems.ProblemPhyXMini0896.lengthInMeters}
  \uses{def:physics:phyx-mini-0896:phyxminiproblems-problemphyxmini0896-lengthquantity}
  Coherent-SI metre readout of a physical length
\end{definition}
\begin{proof}
  This is the declared model definition of length In Meters.
\end{proof}
\begin{definition}[length In Centimeters]
  \label{def:physics:phyx-mini-0896:phyxminiproblems-problemphyxmini0896-lengthincentimeters}
  \lean{PhyXMiniProblems.ProblemPhyXMini0896.lengthInCentimeters}
  \uses{def:physics:phyx-mini-0896:phyxminiproblems-problemphyxmini0896-lengthquantity, def:physics:phyx-mini-0896:phyxminiproblems-problemphyxmini0896-lengthinmeters}
  Centimetre readout used by the two dimension labels in the image
\end{definition}
\begin{proof}
  This is the declared model definition of length In Centimeters.
\end{proof}
\begin{definition}[force Vector In Newtons]
  \label{def:physics:phyx-mini-0896:phyxminiproblems-problemphyxmini0896-forcevectorinnewtons}
  \lean{PhyXMiniProblems.ProblemPhyXMini0896.forceVectorInNewtons}
  \uses{def:physics:phyx-mini-0896:phyxminiproblems-problemphyxmini0896-planarforcequantity}
  Coherent-SI planar-force vector, whose coordinates are in newtons
\end{definition}
\begin{proof}
  This is the declared model definition of force Vector In Newtons.
\end{proof}
\begin{definition}[force Magnitude In Newtons]
  \label{def:physics:phyx-mini-0896:phyxminiproblems-problemphyxmini0896-forcemagnitudeinnewtons}
  \lean{PhyXMiniProblems.ProblemPhyXMini0896.forceMagnitudeInNewtons}
  \uses{def:physics:phyx-mini-0896:phyxminiproblems-problemphyxmini0896-planarforcequantity, def:physics:phyx-mini-0896:phyxminiproblems-problemphyxmini0896-forcevectorinnewtons}
  Euclidean magnitude of a planar physical force, in newtons
\end{definition}
\begin{proof}
  This is the declared model definition of force Magnitude In Newtons.
\end{proof}
\begin{definition}[Charge Site]
  \label{def:physics:phyx-mini-0896:phyxminiproblems-problemphyxmini0896-chargesite}
  \lean{PhyXMiniProblems.ProblemPhyXMini0896.ChargeSite}
  The three occupied corners in the supplied image
\end{definition}
\begin{proof}
  This is the declared model definition of Charge Site.
\end{proof}
\begin{definition}[Source Charge]
  \label{def:physics:phyx-mini-0896:phyxminiproblems-problemphyxmini0896-sourcecharge}
  \lean{PhyXMiniProblems.ProblemPhyXMini0896.SourceCharge}
  The two charges exerting force on the queried lower-right charge
\end{definition}
\begin{proof}
  This is the declared model definition of Source Charge.
\end{proof}
\begin{definition}[source Site]
  \label{def:physics:phyx-mini-0896:phyxminiproblems-problemphyxmini0896-sourcesite}
  \lean{PhyXMiniProblems.ProblemPhyXMini0896.sourceSite}
  \uses{def:physics:phyx-mini-0896:phyxminiproblems-problemphyxmini0896-chargesite, def:physics:phyx-mini-0896:phyxminiproblems-problemphyxmini0896-sourcecharge}
  The site occupied by each source charge
\end{definition}
\begin{proof}
  This is the declared model definition of source Site.
\end{proof}
\begin{definition}[Figure Charge Color]
  \label{def:physics:phyx-mini-0896:phyxminiproblems-problemphyxmini0896-figurechargecolor}
  \lean{PhyXMiniProblems.ProblemPhyXMini0896.FigureChargeColor}
  The red and teal circle fills used in the primary raster
\end{definition}
\begin{proof}
  This is the declared model definition of Figure Charge Color.
\end{proof}
\begin{definition}[Figure Charge Sign]
  \label{def:physics:phyx-mini-0896:phyxminiproblems-problemphyxmini0896-figurechargesign}
  \lean{PhyXMiniProblems.ProblemPhyXMini0896.FigureChargeSign}
  The sign glyph drawn inside each charge marker
\end{definition}
\begin{proof}
  This is the declared model definition of Figure Charge Sign.
\end{proof}
\begin{definition}[expected Charge Color]
  \label{def:physics:phyx-mini-0896:phyxminiproblems-problemphyxmini0896-expectedchargecolor}
  \lean{PhyXMiniProblems.ProblemPhyXMini0896.expectedChargeColor}
  \uses{def:physics:phyx-mini-0896:phyxminiproblems-problemphyxmini0896-chargesite, def:physics:phyx-mini-0896:phyxminiproblems-problemphyxmini0896-figurechargecolor}
  Expected marker colour at each occupied corner
\end{definition}
\begin{proof}
  This is the declared model definition of expected Charge Color.
\end{proof}
\begin{definition}[expected Charge Sign]
  \label{def:physics:phyx-mini-0896:phyxminiproblems-problemphyxmini0896-expectedchargesign}
  \lean{PhyXMiniProblems.ProblemPhyXMini0896.expectedChargeSign}
  \uses{def:physics:phyx-mini-0896:phyxminiproblems-problemphyxmini0896-chargesite, def:physics:phyx-mini-0896:phyxminiproblems-problemphyxmini0896-figurechargesign}
  Expected sign glyph at each occupied corner
\end{definition}
\begin{proof}
  This is the declared model definition of expected Charge Sign.
\end{proof}
\begin{definition}[expected Charge In Nanocoulombs]
  \label{def:physics:phyx-mini-0896:phyxminiproblems-problemphyxmini0896-expectedchargeinnanocoulombs}
  \lean{PhyXMiniProblems.ProblemPhyXMini0896.expectedChargeInNanocoulombs}
  \uses{def:physics:phyx-mini-0896:phyxminiproblems-problemphyxmini0896-chargesite}
  Signed nanocoulomb value printed beside each charge marker
\end{definition}
\begin{proof}
  This is the declared model definition of expected Charge In Nanocoulombs.
\end{proof}
\begin{definition}[Figure Separation]
  \label{def:physics:phyx-mini-0896:phyxminiproblems-problemphyxmini0896-figureseparation}
  \lean{PhyXMiniProblems.ProblemPhyXMini0896.FigureSeparation}
  The two dashed separations carrying numerical labels
\end{definition}
\begin{proof}
  This is the declared model definition of Figure Separation.
\end{proof}
\begin{definition}[Three Point Charge Figure]
  \label{def:physics:phyx-mini-0896:phyxminiproblems-problemphyxmini0896-threepointchargefigure}
  \lean{PhyXMiniProblems.ProblemPhyXMini0896.ThreePointChargeFigure}
  \uses{def:physics:phyx-mini-0896:phyxminiproblems-problemphyxmini0896-chargesite, def:physics:phyx-mini-0896:phyxminiproblems-problemphyxmini0896-figurechargecolor, def:physics:phyx-mini-0896:phyxminiproblems-problemphyxmini0896-figurechargesign, def:physics:phyx-mini-0896:phyxminiproblems-problemphyxmini0896-figureseparation}
  Literal presentation data transcribed from image 896.png
\end{definition}
\begin{proof}
  This is the declared model definition of Three Point Charge Figure.
\end{proof}
\begin{definition}[position Vector In Meters]
  \label{def:physics:phyx-mini-0896:phyxminiproblems-problemphyxmini0896-positionvectorinmeters}
  \lean{PhyXMiniProblems.ProblemPhyXMini0896.positionVectorInMeters}
  Turn a planar Space 2 point into its explicitly metre-valued vector
\end{definition}
\begin{proof}
  This is the declared model definition of position Vector In Meters.
\end{proof}
\begin{definition}[Three Point Charge Setup]
  \label{def:physics:phyx-mini-0896:phyxminiproblems-problemphyxmini0896-threepointchargesetup}
  \lean{PhyXMiniProblems.ProblemPhyXMini0896.ThreePointChargeSetup}
  \uses{def:physics:phyx-mini-0896:phyxminiproblems-problemphyxmini0896-signedchargequantity, def:physics:phyx-mini-0896:phyxminiproblems-problemphyxmini0896-lengthquantity, def:physics:phyx-mini-0896:phyxminiproblems-problemphyxmini0896-planarforcequantity, def:physics:phyx-mini-0896:phyxminiproblems-problemphyxmini0896-chargesite, def:physics:phyx-mini-0896:phyxminiproblems-problemphyxmini0896-sourcecharge, def:physics:phyx-mini-0896:phyxminiproblems-problemphyxmini0896-threepointchargefigure}
  The independent physical objects in the electrostatic setup. The pair-force vectors and resultant force are observables constrained only by the governing laws below; no numerical answer value is built into them
\end{definition}
\begin{proof}
  This is the declared model definition of Three Point Charge Setup.
\end{proof}
\begin{definition}[displacement From Source To Target In Meters]
  \label{def:physics:phyx-mini-0896:phyxminiproblems-problemphyxmini0896-displacementfromsourcetotargetinmeters}
  \lean{PhyXMiniProblems.ProblemPhyXMini0896.displacementFromSourceToTargetInMeters}
  \uses{def:physics:phyx-mini-0896:phyxminiproblems-problemphyxmini0896-sourcecharge, def:physics:phyx-mini-0896:phyxminiproblems-problemphyxmini0896-sourcesite, def:physics:phyx-mini-0896:phyxminiproblems-problemphyxmini0896-positionvectorinmeters, def:physics:phyx-mini-0896:phyxminiproblems-problemphyxmini0896-threepointchargesetup}
  Displacement from a named source to the force target, in metres
\end{definition}
\begin{proof}
  This is the declared model definition of displacement From Source To Target In Meters.
\end{proof}
\begin{definition}[Matches Question Target]
  \label{def:physics:phyx-mini-0896:phyxminiproblems-problemphyxmini0896-matchesquestiontarget}
  \lean{PhyXMiniProblems.ProblemPhyXMini0896.MatchesQuestionTarget}
  \uses{def:physics:phyx-mini-0896:phyxminiproblems-problemphyxmini0896-threepointchargesetup}
  The question asks for the force on the lower-right -10 nC charge
\end{definition}
\begin{proof}
  This is the declared model definition of Matches Question Target.
\end{proof}
\begin{definition}[Matches Supplied Three Charge Figure]
  \label{def:physics:phyx-mini-0896:phyxminiproblems-problemphyxmini0896-matchessuppliedthreechargefigure}
  \lean{PhyXMiniProblems.ProblemPhyXMini0896.MatchesSuppliedThreeChargeFigure}
  \uses{def:physics:phyx-mini-0896:phyxminiproblems-problemphyxmini0896-chargeinnanocoulombs, def:physics:phyx-mini-0896:phyxminiproblems-problemphyxmini0896-lengthinmeters, def:physics:phyx-mini-0896:phyxminiproblems-problemphyxmini0896-lengthincentimeters, def:physics:phyx-mini-0896:phyxminiproblems-problemphyxmini0896-expectedchargecolor, def:physics:phyx-mini-0896:phyxminiproblems-problemphyxmini0896-expectedchargesign, def:physics:phyx-mini-0896:phyxminiproblems-problemphyxmini0896-expectedchargeinnanocoulombs, def:physics:phyx-mini-0896:phyxminiproblems-problemphyxmini0896-positionvectorinmeters, def:physics:phyx-mini-0896:phyxminiproblems-problemphyxmini0896-threepointchargesetup}
  All literal labels and geometric readouts from the primary raster. In particular, the 3.0 cm label belongs to the lower horizontal segment between the +10 nC and -10 nC charges, while the 1.0 cm label belongs to the right vertical segment
\end{definition}
\begin{proof}
  This is the declared model definition of Matches Supplied Three Charge Figure.
\end{proof}
\begin{definition}[Uses Vacuum Coulomb Constant]
  \label{def:physics:phyx-mini-0896:phyxminiproblems-problemphyxmini0896-usesvacuumcoulombconstant}
  \lean{PhyXMiniProblems.ProblemPhyXMini0896.UsesVacuumCoulombConstant}
  \uses{def:physics:phyx-mini-0896:phyxminiproblems-problemphyxmini0896-threepointchargesetup}
  The standard free-space calibration of Coulomb's constant. Its scalar readout has coherent-SI units N m²/C² in the force law below
\end{definition}
\begin{proof}
  This is the declared model definition of Uses Vacuum Coulomb Constant.
\end{proof}
\begin{definition}[Has Physical Three Charge Parameters]
  \label{def:physics:phyx-mini-0896:phyxminiproblems-problemphyxmini0896-hasphysicalthreechargeparameters}
  \lean{PhyXMiniProblems.ProblemPhyXMini0896.HasPhysicalThreeChargeParameters}
  \uses{def:physics:phyx-mini-0896:phyxminiproblems-problemphyxmini0896-lengthinmeters, def:physics:phyx-mini-0896:phyxminiproblems-problemphyxmini0896-threepointchargesetup, def:physics:phyx-mini-0896:phyxminiproblems-problemphyxmini0896-displacementfromsourcetotargetinmeters}
  Positivity and source-target separation for the depicted branch
\end{definition}
\begin{proof}
  This is the declared model definition of Has Physical Three Charge Parameters.
\end{proof}
\begin{definition}[Satisfies Point Charge Coulomb Force Law]
  \label{def:physics:phyx-mini-0896:phyxminiproblems-problemphyxmini0896-satisfiespointchargecoulombforcelaw}
  \lean{PhyXMiniProblems.ProblemPhyXMini0896.SatisfiesPointChargeCoulombForceLaw}
  \uses{def:physics:phyx-mini-0896:phyxminiproblems-problemphyxmini0896-chargeincoulombs, def:physics:phyx-mini-0896:phyxminiproblems-problemphyxmini0896-forcevectorinnewtons, def:physics:phyx-mini-0896:phyxminiproblems-problemphyxmini0896-sourcesite, def:physics:phyx-mini-0896:phyxminiproblems-problemphyxmini0896-threepointchargesetup, def:physics:phyx-mini-0896:phyxminiproblems-problemphyxmini0896-displacementfromsourcetotargetinmeters}
  For source charge q\_s, target charge q\_t, and displacement r from the source to the target, the force on the target is k q\_s q\_t r / ‖r‖³. The signs of the charges therefore determine attraction or repulsion. This relation contains no requested force magnitude
\end{definition}
\begin{proof}
  This is the declared model definition of Satisfies Point Charge Coulomb Force Law.
\end{proof}
\begin{definition}[Satisfies Electrostatic Force Superposition]
  \label{def:physics:phyx-mini-0896:phyxminiproblems-problemphyxmini0896-satisfieselectrostaticforcesuperposition}
  \lean{PhyXMiniProblems.ProblemPhyXMini0896.SatisfiesElectrostaticForceSuperposition}
  \uses{def:physics:phyx-mini-0896:phyxminiproblems-problemphyxmini0896-forcevectorinnewtons, def:physics:phyx-mini-0896:phyxminiproblems-problemphyxmini0896-sourcecharge, def:physics:phyx-mini-0896:phyxminiproblems-problemphyxmini0896-threepointchargesetup}
  The net force on the queried charge is the vector sum of both pair forces
\end{definition}
\begin{proof}
  This is the declared model definition of Satisfies Electrostatic Force Superposition.
\end{proof}
\begin{definition}[Answer Choice]
  \label{def:physics:phyx-mini-0896:phyxminiproblems-problemphyxmini0896-answerchoice}
  \lean{PhyXMiniProblems.ProblemPhyXMini0896.AnswerChoice}
  Labels attached to the four multiple-choice answers
\end{definition}
\begin{proof}
  This is the declared model definition of Answer Choice.
\end{proof}
\begin{definition}[force Magnitude In Newtons]
  \label{def:physics:phyx-mini-0896:phyxminiproblems-problemphyxmini0896-answerchoice-forcemagnitudeinnewtons}
  \lean{PhyXMiniProblems.ProblemPhyXMini0896.AnswerChoice.forceMagnitudeInNewtons}
  \uses{def:physics:phyx-mini-0896:phyxminiproblems-problemphyxmini0896-forcemagnitudeinnewtons, def:physics:phyx-mini-0896:phyxminiproblems-problemphyxmini0896-answerchoice}
  Force magnitude in newtons printed by each answer choice
\end{definition}
\begin{proof}
  This is the declared model definition of force Magnitude In Newtons.
\end{proof}
\begin{definition}[recorded Dataset Answer]
  \label{def:physics:phyx-mini-0896:phyxminiproblems-problemphyxmini0896-recordeddatasetanswer}
  \lean{PhyXMiniProblems.ProblemPhyXMini0896.recordedDatasetAnswer}
  \uses{def:physics:phyx-mini-0896:phyxminiproblems-problemphyxmini0896-answerchoice}
  The answer label recorded in the supplied dataset metadata
\end{definition}
\begin{proof}
  This is the declared model definition of recorded Dataset Answer.
\end{proof}
\begin{definition}[Rounds To Nearest Quantum]
  \label{def:physics:phyx-mini-0896:phyxminiproblems-problemphyxmini0896-roundstonearestquantum}
  \lean{PhyXMiniProblems.ProblemPhyXMini0896.RoundsToNearestQuantum}
  displayed is a nearest-multiple-of-quantum rounding of actual. Requiring the displayed value to lie on the quantum grid distinguishes the two-significant- figure answer 7.3 * 10\textasciicircum{}-3 N from an unrounded decimal nearby
\end{definition}
\begin{proof}
  This is the declared model definition of Rounds To Nearest Quantum.
\end{proof}
\begin{lemma}[net Force Vector eq coulomb Sum]
  \label{lem:physics:phyx-mini-0896:phyxminiproblems-problemphyxmini0896-netforcevector-eq-coulombsum}
  \lean{PhyXMiniProblems.ProblemPhyXMini0896.netForceVector_eq_coulombSum}
  \uses{def:physics:phyx-mini-0896:phyxminiproblems-problemphyxmini0896-chargeincoulombs, def:physics:phyx-mini-0896:phyxminiproblems-problemphyxmini0896-forcevectorinnewtons, def:physics:phyx-mini-0896:phyxminiproblems-problemphyxmini0896-sourcecharge, def:physics:phyx-mini-0896:phyxminiproblems-problemphyxmini0896-sourcesite, def:physics:phyx-mini-0896:phyxminiproblems-problemphyxmini0896-threepointchargesetup, def:physics:phyx-mini-0896:phyxminiproblems-problemphyxmini0896-displacementfromsourcetotargetinmeters, def:physics:phyx-mini-0896:phyxminiproblems-problemphyxmini0896-hasphysicalthreechargeparameters, def:physics:phyx-mini-0896:phyxminiproblems-problemphyxmini0896-satisfiespointchargecoulombforcelaw, def:physics:phyx-mini-0896:phyxminiproblems-problemphyxmini0896-satisfieselectrostaticforcesuperposition}
  Coulomb's law and superposition give the net force as a sum depending only on the two sources, the target charge, their positions, and Coulomb's constant
\end{lemma}
\begin{proof}
  The conclusion follows from the governing relations and figure readouts named in the statement of net Force Vector eq coulomb Sum.
\end{proof}
\begin{lemma}[net Force Magnitude numerical Bounds]
  \label{lem:physics:phyx-mini-0896:phyxminiproblems-problemphyxmini0896-netforcemagnitude-numericalbounds}
  \lean{PhyXMiniProblems.ProblemPhyXMini0896.netForceMagnitude_numericalBounds}
  \uses{def:physics:phyx-mini-0896:phyxminiproblems-problemphyxmini0896-forcemagnitudeinnewtons, def:physics:phyx-mini-0896:phyxminiproblems-problemphyxmini0896-threepointchargesetup, def:physics:phyx-mini-0896:phyxminiproblems-problemphyxmini0896-matchesquestiontarget, def:physics:phyx-mini-0896:phyxminiproblems-problemphyxmini0896-matchessuppliedthreechargefigure, def:physics:phyx-mini-0896:phyxminiproblems-problemphyxmini0896-usesvacuumcoulombconstant, def:physics:phyx-mini-0896:phyxminiproblems-problemphyxmini0896-hasphysicalthreechargeparameters, def:physics:phyx-mini-0896:phyxminiproblems-problemphyxmini0896-satisfiespointchargecoulombforcelaw, def:physics:phyx-mini-0896:phyxminiproblems-problemphyxmini0896-satisfieselectrostaticforcesuperposition}
  For the charge and distance readouts in the image, the unrounded magnitude is approximately 7.2591 * 10\textasciicircum{}-3 N. These bounds are exactly the interval that rounds to 7.3 * 10\textasciicircum{}-3 N at resolution 10\textasciicircum{}-4 N
\end{lemma}
\begin{proof}
  The conclusion follows from the governing relations and figure readouts named in the statement of net Force Magnitude numerical Bounds.
\end{proof}
% --- Archon declaration topology end ---
% --- Archon physics formalization source end ---
